\hypertarget{cont10:paper-ix-simultaneous-coefficient-and-parameter-deformation}{%
\subsection{Paper IX: simultaneous coefficient and parameter
deformation}\label{cont10:paper-ix-simultaneous-coefficient-and-parameter-deformation}}

\textbf{Author: Komorebi P. --- 8 September 2026.}

This block combines coefficient perturbation with prescribed parameter
motion. The original solved Jacobian controls existence, local
uniqueness, displacement and the first-order error. The four new
mathematical modules compile. The fresh verification report and
independent audit record the authoritative release status and exact
source revisions.

\hypertarget{cont10:the-theorem}{%
\subsubsection{The theorem}\label{cont10:the-theorem}}

Let \(p\) be any prime, \(R=\mathbb Z_p\), and \(Q=\mathbb Q_p\). Let \[
F,H\in R[Y_1,\ldots,Y_e,T_1,\ldots,T_k]^e,
\qquad a=(y_0,t_0)\in R^{e+k},\qquad F(a)=0.
\] Here \(e,k\ge0\) are arbitrary finite dimensions. There are as many
solved coordinates as equations. Define the actual derivative blocks at
the original root: \[
J=D_YF(a)\in M_e(R),\qquad K=D_TF(a)\in M_{e\times k}(R).
\] Suppose, for natural numbers \(s,M\), that \[
J\in\operatorname{GL}_e(Q),\qquad p^sJ^{-1}\in M_e(R),\qquad M>2s.
\] For every \(v\in R^k\), set \[
\tau=t_0+p^Mv,\qquad G=F+p^MH.
\] Then there is exactly one \(y\in y_0+p^{s+1}R^e\) with \[
G(y,\tau)=0.
\] It satisfies \[
y-y_0\in p^{M-s}R^e
\] and there is an integral vector \(\varepsilon\in p^{2M-3s}R^e\) such
that, in \(Q^e\), \[
\boxed{\displaystyle
 y-y_0=-p^MJ^{-1}\bigl(H(a)+Kv\bigr)+\varepsilon.
}
\] The same original exponent \(s\) appears throughout. No inverse or
rank assumption on the perturbed Jacobian is imposed. The statement
includes arbitrary integral p-adic coefficients and centers, the dyadic
prime, and zero dimensions.

The full-point theorem constructs \(b=(y,\tau)\) and also proves \[
b-a\in p^{M-s}R^{e+k}\cap p^{s+1}R^{e+k}.
\] Parameter equality \(b_T=\tau\) is exact. Uniqueness concerns every
root in the stated solved-coordinate ball with these parameters, not
only roots obtained by the construction. No global uniqueness is
claimed.

\hypertarget{cont10:actual-augmented-system}{%
\subsubsection{Actual augmented system}\label{cont10:actual-augmented-system}}

Define polynomial systems with \(e+k\) equations and variables by \[
\mathcal F(Y,T)=\bigl(F(Y,T),T-t_0\bigr),\qquad
\mathcal H(Y,T)=\bigl(H(Y,T),-v\bigr).
\] Both are actual polynomials over \(R\). Direct evaluation proves \[
\mathcal F(a)=0,
\qquad
(\mathcal F+p^M\mathcal H)(b)=0
\iff F(b)+p^MH(b)=0\ \text{and}\ b_T=t_0+p^Mv.
\] Actual polynomial partial differentiation gives \[
D\mathcal F(a)=B=\begin{pmatrix}J&K\\0&I_k\end{pmatrix}.
\] These identities are proved from evaluation and partial derivatives,
rather than supplied as additional hypotheses.

Choose the integral matrix \(A=p^sJ^{-1}\). The existing
inverse-certificate equivalence gives \[
AJ=JA=p^sI_e.
\] Put \(c=p^s\) and define an integral block matrix \[
C=\begin{pmatrix}A&-AK\\0&cI_k\end{pmatrix}.
\] Then \[
CB=\begin{pmatrix}AJ&AK-AK\\0&cI_k\end{pmatrix}=cI_{e+k},
\] \[
BC=\begin{pmatrix}JA&-JAK+cK\\0&cI_k\end{pmatrix}=cI_{e+k}.
\] The second equality uses \(JA=cI_e\). The generic block proof works
over any commutative ring. The p-adic specialization also explicitly
proves that \(B\) is invertible over \(Q\) and that \(p^sB^{-1}\) is
integral. There is no denominator loss from appending the parameter
equations.

\hypertarget{cont10:construction-error-and-uniqueness}{%
\subsubsection{Construction, error and
uniqueness}\label{cont10:construction-error-and-uniqueness}}

Apply the previously checked weighted coefficient-perturbation
construction to \(\mathcal F,\mathcal H,a,C\). Its hypotheses follow
from the actual augmented identities and the two-sided certificate. It
constructs \(b\), proves the full displacement bound \(p^{M-s}\), and
proves uniqueness throughout \(a+p^{s+1}R^{e+k}\). The augmented root
equivalence gives the required original equations and exact parameter
values.

The checked integral first-order estimate gives \[
b-a+p^{M-s}C\mathcal H(a)\in p^{2M-3s}R^{e+k}.
\] The decisive block action is \[
C\binom{H(a)}{-v}
=
\binom{A\bigl(H(a)+Kv\bigr)}{-p^sv}.
\] Projecting the solved coordinates therefore gives \[
y-y_0+p^{M-s}A\bigl(H(a)+Kv\bigr)
\in p^{2M-3s}R^e.
\] The existing rational-conversion theorem turns \(p^{M-s}A\) into
\(p^MJ^{-1}\) and retains an integral error vector. This proves the
displayed first-order formula with the original derivative blocks.

For any competing solved-coordinate root \(z\in y_0+p^{s+1}R^e\), the
joined point \((z,\tau)\) lies in the augmented uniqueness ball:
\(M>2s\) implies \(M\ge s+1\), so \(\tau-t_0=p^Mv\in p^{s+1}R^k\).
Full-point uniqueness gives \((z,\tau)=b\), hence \(z=y\). This is the
additional comparison needed to obtain uniqueness in exactly the ball
stated in the solved-coordinate theorem.

\hypertarget{cont10:checked-controls-and-precision-boundary}{%
\subsubsection{Checked controls and precision
boundary}\label{cont10:checked-controls-and-precision-boundary}}

For \(p=2\), \(s=0\), \(F(Y,T)=Y+T\), \(H=1\), \(a=0\), and \(v=1\),
both original blocks satisfy \(J=K=1\). At \(M=1\), the exact parameter
is \(\tau=2\), and the unique solved value is \(y=-4\). The combined
correction is \(-2(1+1)=-4\), with zero error.

The sign controls use \(M=2\): the actual augmented root is
\((y,\tau)=(-8,4)\), and the required error ideal is \(16\mathbb Z_2\).
Omitting \(Kv\) predicts \(-4\), with error \(-4\notin16\mathbb Z_2\).
Reversing its sign predicts \(0\), with error \(-8\notin16\mathbb Z_2\).
Both failures are checked. At binary depth \(M=1\), the reversed-sign
difference would lie in the allowed error ideal, so that shallower
example alone would not distinguish the sign. This control-design
limitation was identified and corrected before the release gate.

For parameter-only motion, use \[
F(Y,T)=Y+Y^2+T,\quad H=0,\quad a=0,\quad v=1,\quad p=2,\ s=0,\ M=1.
\] A checked transfer from the prior scalar quadratic root constructs an
even \(y\) with \[
y+y^2+2=0,\qquad 4\mid y+2,\qquad 8\nmid y+2.
\] Thus the error exponent \(2M-3s=2\) is attained even when only
parameters move. No uniform extra error digit can be added to the
general theorem. This example establishes that limitation; it does not
assert optimality at every choice of \(p,s,M\).

The inherited controls continue to record the necessity of the strict
weighted threshold and other earlier scope boundaries. General
rationally rank-deficient systems are outside the inverse hypothesis.
